% !TeX root = ../Masterarbeit.tex

\documentclass[Masterarbeit.tex]{subfile}

\begin{document}
	
	\chapter{Discussion} \label{discussion}
	
	The results produced in this thesis, the applied methods and the executed course of action are all based on assumptions. These assumptions and the neural network's limitations have to be discussed first before any final conclusions can be drawn in regard to its performance and potential. Additionally, it can be examined in how far these limitations are specific to this work and can be worked on in further research outside of this thesis. \\
	\begin{figure}[H]
		\centering
		\includegraphics[width=1.0\linewidth]{plots/summary_all_estimates}
		\caption{Comparison of all utilized SPG OHC estimates: EnKF assimilation ensemble mean with its spread (red), the indirect neural network reconstructions with their ensemble spread (blue), the direct neural network reconstructions with their ensemble spread (green) and the IAP (orange) and EN4 (purple) SPG OHC timeseries.}
		\label{fig:summaryallestimates}
	\end{figure}
		
	The neural networks are only trained on temperature profiles produced by the EnKF assimilations described in section \ref{data:assim}. Therefore they are the neural network's main point of reference and the basis for its evaluation. Thus, the entire analysis bases on the assumption that these assimilations provide a close to accurate description of the physical reality during the Argo era. In the training process, the neural network learns to reproduce the assimilations spatial patterns and overall trends. In turn, that means that the neural network is limited by its training data's uncertainties and errors. While it can transfer these patterns to times with lower observational densities, it can never outperform but at most reproduce the assimilation in times with the equal or more observations compared to the training period. Additionally, it also reproduces any uncertainties or errors that are contained in its training data as well as the spatial patterns that are produced by the assimilations specific set of parameters and inputs. Figure \ref{fig:summaryallestimates} shows that estimates of SPG OHC can differ significantly and that the neural network is clearly biased towards reproducing its training data. A way to counter this bias and at the same time better estimate the uncertainties of the assimilation during the training period could be the implementation of multi model training data for the neural network. By introducing the training data from different but comparable assimilations, it could help identifying uncertainties due to model biases and uncertainties. \\
	In addition to the limitation due to its training data, the neural network reconstructions have also shown do be dependent on the time period from which their training data is taken. This work is based on the assumption that the physical patterns learned by the neural network during the Argo era are directly applicable to the reconstruction during the pre-Argo era. However, natural variability with periods longer than the training time as well as changes due to anthropogenic influence can impact the reconstructed variables significantly. Thus, it has to be further examined if a training period of 17 years (the Argo era up to 2020) is long enough for the neural network to capture the impact of different modes of low frequency natural variability and anthropogenic influence. Similarly, it also remains to be examined how far into the future the neural networks reconstructions can be applied with a training period of 2004 -- 2020. Therefore, the training period has to be extended at some point when reconstructing newer timesteps, especially with rapid changes in regional patterns, such as Arctic heat waves (cite Armineh). Furthermore, when examining how far into the past the neural network reconstructions are valid, the question also arises how many points of observation are required by the neural network. The results of this work show that up to 1958 the neural network is able to reproduce the assimilation. However, before 1958 the observational density in the NA decreases even further. Which amount of observations is necessary for the neural network in order to still produce significant results, should be examined in further research. \\	
	Another key limitation of this work, which has already been indicated, is the utility of the reconstruction created in this thesis. All analysis steps have focused on directly comparing the different OHC estimates without examining its further uses. In current research, assimilations are mostly used to initialize climate hindcasts \cite{brune_assimilation_2015}. If the reconstructions produced by partial convolutional U-nets are also able to provide hindcasts with meaningful initialization has yet to be examined. Therefore, replacing or even only assisting EnKF assimilations with neural network reconstructions can only be seriously considered after its utility for initialization has been proven. \\
	On a more technical level, section \ref{uncertainty} discussed some of the largest uncertainties contained in the neural network and its training process. Nonetheless, these uncertainties are also based on the assumption that the calculated metrics accurately capture the neural network's ability to reproduce the assimilation. Section \ref{analysis_preargo} has already shown that the loss parameters have significant shortcomings in measuring the neural network's performance. Therefore, it can be assumed that for example the choice of stopping point as a result of the examination of the validation loss still does not capture the full extend of the contained uncertainty. This uncertainty could be limited by further hyperparameter training of the loss parameters or the learning rate. But not only the hyperparameter training could be improved. The neural network could also be assisted by the implementation of other additional functions. On the one hand, machine learning techniques, such as long-short-term-memory (LSTM) or attention modules could be provided. Especially LSTM modules could prove to be helpful by providing information in the time dimension to the reconstruction. On the other hand the neural network could be improved by adding additional physical information, such as additional variables in the different input channels. In the case of the reconstruction of temperature profiles, the addition of variables, such as salinity or 2m-air temperature could be considered. \\
	In general, the application of this partial convolutional U-net to the reconstruction of other variables other than temperature profiles could examined. Similarly, the reconstruction of temperature profiles in other regions than the NA could be tested. It would be especially interesting to explore in how far the neural network reconstruction still represents a valid alternative to the EnKF assimilations in regions with higher scarcity of observations. \\
	Finally, all estimates of required computational resources in section \ref{resources} have been made for just this 128x128 grid of the NA. Thus, it could be examined how the differences in computational requirements for both the neural network and EnKF assimilation change with a higher resolution, larger ensembles or larger areas. This could lead to a better estimate of how much time and computational resources the work with machine learning can save. Furthermore, it could explored if neural networks could push resolution of assimilations beyond the EnKF assimilation's current limitations. \\
	

		
		
		

		
		
		
		
		
		
		
	
	
	
	
	
	
	
	
	
	
	
	
	
	
	
\end{document}